% =====================================================================
% Section 2 body fragment — Shock-Geometry Pre-Pass Architecture
% Parent jsr_paper.tex declares \section{Shock-Geometry Pre-Pass Architecture}
% Numbers sourced from _DRAFTING_BRIEF.md (authoritative). [F3] honest framing applied.
% =====================================================================

This section presents the architectural contribution that distinguishes the
present work from prior Barrowman-pipeline open-source simulators. The
shock-geometry pre-pass is a once-per-timestep nose-to-tail surface march that
publishes the post-shock local Mach number, static pressure, and static
temperature at every component station before any downstream drag or stability
calculator runs. Subsequent sections consume those local conditions through a
stable query interface. The architectural value of the pre-pass lies in local
post-shock fidelity and in providing a single seam through which all of the
supersonic submodels of Sections~3--5 are driven; as quantified in
Section~\ref{subsec:prepass-ablation}, its contribution to \emph{integrated}
apogee is small, because apogee integrates a trajectory whose drag impulse is
dominated by lower-Mach flight.

\subsection{Motivation}

The freestream Mach number governs aerodynamic-coefficient behavior only when
the flow over each component is unprocessed by upstream shocks and expansions.
On a slender finned rocket at moderate supersonic Mach this assumption breaks at
the first body shoulder: the conical or ogival nose shock processes the flow
before it reaches the cylindrical afterbody, the fins, and the base. A
\SI{15}{\degree} cone at freestream Mach~2.5, for instance, leaves the
post-shock surface Mach near 2.16---a fourteen-percent reduction. Because the
Prandtl--Glauert compressibility factor $\beta = \sqrt{|M^2 - 1|}$ is nonlinear,
that fourteen-percent Mach reduction translates to roughly an eighteen-percent
change in the fin lift-curve-slope coefficient $K_1 = 2/\beta$. Similar effects
propagate through fin wave drag, fin--body interference, and base-region inflow
conditions. Evaluating every downstream coefficient at the freestream Mach
therefore introduces a systematic bias in the \emph{local} aerodynamic loads,
even where its effect on the integrated trajectory is modest.

Closed-source missile aerodynamic codes such as Missile~DATCOM and the
Aeroprediction lineage~\cite{sooy2005,moore2001,mooremcinvillehymer2002} handle
post-shock local conditions internally as part of their coefficient generators
but do not expose them as a shared architectural seam. Open-source rocket
simulators do not address the local-flow question at all: OpenRocket and
RocketPy evaluate every component at the freestream
Mach~\cite{niskanen2009,rocketpy2021}. The present pre-pass makes the
post-shock state an explicit, queryable object shared by every downstream
calculator within an open six-degree-of-freedom trajectory simulator.

\subsection{Data flow}
\label{subsec:prepass-dataflow}

The pre-pass executes once per call to the top-level aerodynamic-forces routine,
immediately after the \texttt{FlightConditions} object is populated and before
any per-component calculator runs. It produces a \texttt{ShockGeometry} object
indexed by axial station; each downstream calculator queries that object for its
local post-shock state. The data path is summarized below:

\begin{quote}
\ttfamily\small
function getAerodynamicForces(rocket, fc):\\
\hspace*{1.5em} shock\_geom = ShockGeometry.compute(rocket, fc)\\
\hspace*{1.5em} forces = \{\}\\
\hspace*{1.5em} for component in rocket.components():\\
\hspace*{3.0em} local = shock\_geom.queryAt(component.axialStation())\\
\hspace*{3.0em} forces[component] = calc(component, fc, local)\\
\hspace*{1.5em} return aggregate(forces)
\end{quote}

The pre-pass is configuration-aware: it invalidates and recomputes on geometric
changes (staging, fairing separation, side-booster jettison) but is otherwise
pure with respect to \texttt{FlightConditions}. Below the sonic line the surface
march short-circuits to a pass-through that returns freestream conditions at
every station; the measured subsonic overhead is below ten microseconds per call
on the reference hardware (single-threaded Intel Core~i7), which is below the
measurement floor of the simulator's per-timestep budget. The full
input/output contract is given in Table~\ref{tab:prepass-contract}.

\begin{table*}[t]
\caption{\texttt{ShockGeometry} calculator input/output contract. Subsonic
column gives the pass-through value returned below $M_\infty = 1.0$.}
\label{tab:prepass-contract}
\centering
\begin{tabular}{@{}llll@{}}
\toprule
Field & Direction & Description & Subsonic value \\
\midrule
\texttt{Rocket} configuration & in & Component tree, axial stations, half-angles, shoulder/boat-tail flags & passed through \\
\texttt{FlightConditions.machNumber} & in & Freestream Mach number $M_\infty$ & passed through \\
\texttt{FlightConditions.aoa} & in & Vehicle angle of attack (windward-side bias when present) & unused \\
\texttt{localMach[station]} & out & Post-shock local Mach at each axial station & returns $M_\infty$ \\
\texttt{localPressure[station]} & out & Post-shock static-pressure ratio $p/p_\infty$ & returns 1.0 \\
\texttt{localTemperature[station]} & out & Post-shock static-temperature ratio $T/T_\infty$ & returns 1.0 \\
\texttt{shockStandoff} & out & Detached-shock standoff distance for blunt noses & $\infty$ (no shock) \\
\texttt{activationBlend} & out & $C^1$ blend weight $w \in [0,1]$ through $M = 1.0$--$1.1$ & returns 0.0 \\
\bottomrule
\end{tabular}
\end{table*}

Figure~\ref{fig:prepass-block} shows the data-flow block diagram.

\begin{figure}[t]
\centering
\includegraphics[width=\columnwidth]{shockgeometry_block_diagram}
\caption{Shock-geometry pre-pass data-flow block diagram. Once per timestep the
pre-pass surface-marches the airframe and publishes the local post-shock Mach,
pressure, and temperature; each component calculator queries the local state at
its axial station before computing forces and moments.}
\label{fig:prepass-block}
\end{figure}

\subsection{Surface marching}
\label{subsec:prepass-march}

The march initializes at the nose tip. For a sharp conical nose the inflow is
processed through a Taylor--Maccoll cone solution~\cite{taylormaccoll1933}
evaluated at the local tip half-angle, returning the cone-surface Mach and
pressure as the seed state. For an ogival nose the tip slope is read as an
instantaneous cone half-angle, and the Taylor--Maccoll seed is followed by a
Prandtl--Meyer expansion at each subsequent strip increment. For a blunt nose a
detached normal-shock solution seeds the post-shock state, and the surface march
resumes around the shoulder.

Body-station traversal walks downstream along the axisymmetric body, applying at
each surface element either a Prandtl--Meyer expansion (for a contracting turning
angle, where the surface inclination decreases relative to the freestream) or an
oblique-shock turn (for an expanding turning angle, where the surface inclination
increases). The branch is selected by the sign of the turning angle relative to
the prior surface tangent; turning angles below a numerical threshold
($10^{-9}~\mathrm{rad}$) are skipped to preserve the isentropic state.
Boat-tails and aft-body contractions are handled by the expansion branch;
shoulders, fin-leading-edge regions, and forward conical transitions invoke the
oblique-shock branch.

Pressure and temperature ratios accumulate multiplicatively along the march;
Mach is updated through the Prandtl--Meyer function~\cite{naca1135} for
expansions and through the post-shock $\theta$--$\beta$--$M$ relation for
compressions. The iterative solvers---a Taylor--Maccoll Runge--Kutta-4
integration with 500 angular steps, an oblique-shock bisection on the shock
angle $\beta$, and a Prandtl--Meyer Newton--Raphson inversion---converge to a
$10^{-12}$ relative tolerance.

\subsection{Local-condition query interface}
\label{subsec:prepass-query}

Downstream calculators receive the pre-pass output through a stable query
interface keyed on axial station. The body calculator reads \texttt{localMach}
and \texttt{localPressure} at each body-element centroid for the
shock-expansion strip integral and the crossflow body lift. The fin-set
calculator reads \texttt{localMach} at the fin axial station and uses it in the
DATCOM $K_1/K_2/K_3$ fin normal-force computation and in the DATCOM~4.1.5.1 fin
wave-drag computation. The fin--body interference factors and the
shock-wave/boundary-layer-interaction chord-reduction model draw from the same
interface, and the base-drag calculator reads the post-shock local pressure for
the base dynamic-pressure rescaling. A single seam thus supplies the local flow
state to every supersonic submodel of Sections~3--5.

\subsection{Numerical guards and $C^1$ activation}
\label{subsec:prepass-activation}

The pre-pass is gated by a continuous activation weight $w(M_\infty)$ that ramps
from zero at $M_\infty = 1.0$ to one at $M_\infty = 1.1$ via a smoothstep,
\begin{equation}
w(M_\infty) = 3t^2 - 2t^3, \qquad t = \frac{M_\infty - 1.0}{0.1},
\quad t \in [0,1].
\label{eq:smoothstep}
\end{equation}
Below $M_\infty = 1.0$ the pre-pass returns freestream conditions everywhere;
above $M_\infty = 1.1$ it returns the full post-shock state; through the
activation band each component's local state is the smoothstep blend of the two.
The activation band is $C^1$-continuous in $M_\infty$, which preserves the
smoothness of the assembled drag and stability coefficients required by the
fixed-step Runge--Kutta-4 trajectory integrator.

Additional guards include a minimum surface-element length below which the
expansion/oblique-shock branch is suppressed (avoiding accumulated round-off on
geometries with very fine longitudinal discretization); a hard fallback to
freestream conditions when the computed local Mach descends below~1.0 due to
repeated expansion (preserving physical consistency); and bounds-checking on the
post-shock $\theta$--$\beta$--$M$ solution to ensure $M_2 > 0$ for all admissible
inputs.

\subsection{Verification against analytical references}
\label{subsec:prepass-verification}

The pre-pass was verified against analytical references at six combinations of
freestream Mach and cone half-angle spanning $M_\infty = 1.5$, 2.5, and 5.0 at
$\theta_c = \SI{10}{\degree}$ and \SI{20}{\degree}. Cone-surface Mach reproduced
the Taylor--Maccoll exact solution to $0.00\%$ in all six cases, the residual
falling below the eight-digit print precision of the diagnostic; the underlying
Taylor--Maccoll shock-angle solver matches the tabulated cone solution to within
$0.825\%$ (against a $1\%$ acceptance gate). Shoulder-expansion turning
reproduced the Prandtl--Meyer relation~\cite{naca1135} to better than
$4 \times 10^{-11}\%$ at all tested turning angles between \SI{1}{\degree} and
\SI{25}{\degree}, and the oblique-shock branch reproduced the closed-form turn
relation to a maximum error of $0.021\%$. Figure~\ref{fig:prepass-mach} plots the
surface-Mach march for the six reference cases against the analytical envelope.
This verification level is consistent with the AIAA Editorial Policy on Numerical
and Experimental Accuracy~\cite{aiaa_numerical_policy}, which directs engineering
codes to anchor against analytical or tabulated solutions where available; the
pre-pass is reported as an externally benchmarked subsystem in the Section~6
roll-up.

\begin{figure}[t]
\centering
\includegraphics[width=\columnwidth]{shockgeometry_surface_mach_verification}
\caption{Shock-geometry surface-Mach verification against Taylor--Maccoll cone
flow and Prandtl--Meyer expansion at six $(M_\infty, \theta_c)$ reference points
spanning $M_\infty = 1.5$--$5.0$ and $\theta_c = \SI{10}{\degree}$--%
\SI{20}{\degree}. Cone-surface Mach is reproduced to $0.00\%$ and
shoulder-expansion Mach to better than $4 \times 10^{-11}\%$.}
\label{fig:prepass-mach}
\end{figure}

\subsection{Contribution to integrated apogee: ablation}
\label{subsec:prepass-ablation}

The architectural prominence of the pre-pass should not be confused with a large
effect on integrated apogee. To quantify its contribution honestly, each
supersonic mechanism was disabled in isolation and the resulting change in
predicted apogee was measured across the 24 single-stage corpus flights.
Disabling the shock-geometry pre-pass changes the integrated apogee error by a
mean magnitude of only $0.15$ percentage points, with a maximum of $3.6$
percentage points on a single Mach-2.46 flight (FMJ~BR6). By comparison, the
externally calibrated finned-base drag augmentation is the dominant apogee
driver, with a mean magnitude of $8.10$ percentage points and a maximum of
$39.5$ percentage points; Van~Driest~II skin friction contributes $0.87$
percentage points (mean) and matters chiefly at high Mach, and the DATCOM fin
wave-drag term contributes $0.39$ percentage points.

The small integrated-apogee effect has a straightforward cause: apogee is a
drag-impulse integral over a trajectory whose drag is dominated by lower-Mach
flight, and the pre-pass is inert below $M_\infty = 1.0$ by construction
(Section~\ref{subsec:prepass-activation}). The pre-pass is therefore \emph{not}
presented as a gross-apogee improvement. Its value is twofold: it supplies
correct \emph{local} post-shock conditions for fin loads and static stability---%
verified to the analytical references of
Section~\ref{subsec:prepass-verification}, where freestream evaluation would bias
the local coefficients as shown in the motivation above---and it is the
architectural seam through which the supersonic drag and stability models of
Sections~3--5 are driven. The integrated-apogee accuracy reported in Section~8
rests primarily on those downstream models and on the externally calibrated
finned-base term, not on the pre-pass in isolation.
